\documentclass[11pt,a4paper]{article}
\usepackage{amsmath,amssymb,amsthm}
\usepackage{geometry}
\usepackage{hyperref}
\usepackage{booktabs}
\usepackage{parskip}
\usepackage{xcolor}
\usepackage{graphicx}
\usepackage{multirow}

\geometry{margin=2.5cm}

\newcommand{\E}{\mathbb{E}}
\newcommand{\N}{\mathcal{N}}
\newcommand{\Lcal}{\mathcal{L}}
\newcommand{\R}{\mathbb{R}}
\newcommand{\code}[1]{\texttt{#1}}
\newcommand{\todo}[1]{\textcolor{red}{[TODO: #1]}}

\title{%
  Can a Fisher/Cram\'er-Rao Bound Predict $\theta$'s Empirical Recovery
  Scatter?\\
  \large Full-covariance CRB vs.\ real pixel-likelihood fits, aberrated
  wavefront
}
\author{}
\date{}

\begin{document}
\maketitle
\tableofcontents
\newpage

%%============================================================
\section{Motivation}
\label{sec:motivation}
%%============================================================

Looking at 2D residual plots of $\hat\theta_{\rm best}-\theta_{\rm true}$
for the \code{confocal\_random\_zernike} (aberrated-wavefront) dataset, the
scatter looks Gaussian. That raises the obvious question: can a
Fisher-information / Cram\'er-Rao bound predict that spread analytically,
rather than only measuring it empirically?

This was tried once before and got stuck. From
\code{notes/joint\_profile\_inference.tex} (a sibling investigation,
\S2026-07-20): a Laplace approximation using the \emph{diagonal} of a local
finite-difference Hessian at $\theta$'s MAP was found \textbf{66--220$\times$
tighter} than the true empirical scatter across 50 shots. That work diagnosed
why, and flagged two things that were never actually done:

\begin{enumerate}
  \item $\theta$'s failure mode is a \emph{ridge} -- a continuous,
    weakly-identified direction (near-zero Fisher eigenvalue), not discrete
    multimodality like $\phi$'s. A ridge is compatible with local Fisher/CRB
    validity \emph{if} it is close to linear/quadratic near the truth -- but
    whether it actually is was flagged as unchecked (Future Work item 0a in
    that file) and never verified.
  \item The failed attempt used a \textbf{diagonal-only} Hessian, discarding
    exactly the off-diagonal correlations a ridge produces. \code{crb\_signal.py}
    already has \code{build\_joint\_fisher()} (full $17\times17$ joint
    Fisher over $(\phi,\theta_{z0},\theta_{z100})$, pixel-likelihood mode,
    all cross terms) and \code{full\_per\_shot\_covariance()} (inverts the
    \emph{full} matrix) -- but the latter was dead code, never called from
    that module's own driver. Nobody had tried the full-covariance version.
\end{enumerate}

This note finishes that specific, previously-flagged, previously-skipped
check, reusing the existing (but previously unwired) machinery rather than
re-deriving anything. Scope, by explicit direction: \textbf{one dataset}
(the aberrated \code{confocal\_random\_zernike} case), matching CRB to
empirical data. A Gaussian-beam-vs-aberrated comparison, and the downstream
connection to the $\beta$ (signal) CRB, are deliberately deferred.

%%============================================================
\section{Setup}
\label{sec:setup}
%%============================================================

Dataset: \code{data/R20\_N200\_A1000000\_..\_psmapconfocal\_random\_zernike}
(20 runs $\times$ 200 shots, $10^6$ atoms/shot, aberrated wavefront via
\code{phase\_space\_grids.py}'s interpolated-PSMAP path). Cloud prior:
independent Gaussian on each of the 8 kinematic components (position,
velocity, position spread, velocity spread, in $x$ and $y$), $\sigma=10\,\mu$m
(or the matching velocity unit) for all 8. Evaluator geometry:
\textbf{exactly} what \code{non\_phase\_shear/analysis/generate\_kinematic\_estimates.py}'s
\code{best} (pixel-likelihood) fit uses -- \code{bins=32}, tight half-range
$1.54\times10^{-3}$\,m, \code{pixel\_n\_gh=8} -- so the CRB is evaluated on
the identical pixel/bin configuration that produced the empirical numbers it
is compared against. Using a different geometry would make the comparison
apples-to-oranges before it starts.

All CRB quantities below are evaluated at the \emph{true} $(\phi,\theta_{z0},
\theta_{z100})$ for each shot -- the standard Cram\'er-Rao convention, and
also convenient here since it needs no extra optimizer run (the empirical
$\hat\theta_{\rm best}$ values already come from a separate, already-run
fit).

Implementation: \code{python-scripts/theta\_crb\_vs\_empirical.py} (new),
importing \code{crb\_signal.py}'s Fisher-matrix code rather than duplicating
it.

%%============================================================
\section{Step 1: Is the ridge direction actually linear/quadratic?}
\label{sec:ridge}
%%============================================================

For 10 real shots (\code{run\_000}), built the full $17\times17$ joint
Fisher matrix $H$ at the true $(\phi,\theta_{z0},\theta_{z100})$, found its
smallest eigenvalue $\lambda_{\min}$ and eigenvector $v$ (the ridge
direction; $\sigma_{\rm ridge}=1/\sqrt{\lambda_{\min}}$), and profiled the
\emph{real} Poisson NLL along $v$ at $k\sigma_{\rm ridge}$ for
$k\in\{0.5,1,2,4,8\}$, comparing to the Fisher-implied quadratic
$\tfrac12\lambda_{\min}(k\sigma_{\rm ridge})^2$.

\textbf{Result: not cleanly quadratic.} The ratio (actual $\Delta$NLL /
quadratic $\Delta$NLL) at $k=0.5$--$1$ (the scale that actually matters --
real fluctuations live near $1\sigma$, not $8\sigma$) ranged from
\textbf{$-2.9$ to $+4.3$} across the 10 shots -- including several shots
where the sign itself flips (the true NLL is briefly \emph{lower}, not
higher, than at the anchor point, in the direction the Fisher matrix calls
"increasing"). The ratio does converge toward 1 at larger $k$ (e.g.\ $0.80$--$1.23$
at $k=8$), consistent with the quadratic approximation eventually dominating
far from the anchor, but that is not the regime real shot noise samples.

This is a genuine, unresolved finding, not swept under the rug: the single
worst-conditioned direction is \emph{not} well-approximated by the local
quadratic at the scale that matters. And yet -- see \S\ref{sec:comparison}
-- the \emph{aggregate}, full 8-dimensional marginal covariance built from
the same local Fisher matrix matches the empirical scatter well. The most
likely reconciliation: $v$ (found by minimizing over the full 17-dimensional
$(\phi,\theta_{z0},\theta_{z100})$ space) is a mixed direction, not aligned
with any single physical $\theta$ component's own marginal-variance
direction, so a $2$--$4\times$ local mis-estimate of curvature along $v$
does not necessarily translate into a comparable error in any individual
component's marginal variance after the Schur-complement/inversion mixes
all 17 directions together. This is a hypothesis, not something checked
here -- flagged in \S\ref{sec:open} as follow-up.

%%============================================================
\section{Step 2+3: Full-covariance CRB vs.\ empirical RMSE}
\label{sec:comparison}
%%============================================================

For the same 10 shots, computed \code{full\_per\_shot\_covariance(H)} and
extracted $\sigma_{\theta_{z0}}$, $\sigma_{\theta_{z100}}$ (marginal std per
component), aggregated as root-mean-square across shots. Compared against
the empirical RMSE of $\hat\theta_{\rm best}-\theta_{\rm true}$, pooled over
z0 and z100, from the 13 completed runs of the real
\code{confocal\_random\_zernike} job available at the time of this check
(5200 shot/z-slice points).

\textbf{One data-quality issue found and handled}: the raw empirical RMSE
for $\sigma_{vy}$ was $2.33\times10^{-6}$, a $3.46\times$ mismatch against
the CRB's $6.75\times10^{-7}$ -- while every other component matched to a
few percent. Investigated directly rather than accepted: of 5200 points, a
single shot has a $\sigma_{vy}$ residual of $1.6\times10^{-4}$ ($161\,\mu$m)
-- roughly $70\times$ larger than the next-largest residual in that
component and wildly inconsistent with the rest of the distribution (median
$|{\rm residual}|=4.4\times10^{-7}$, matching $\sigma_{vx}$'s median almost
exactly). RMSE is quadratic in residuals, so this one likely optimizer
failure (an L-BFGS-B run landing in a bad basin for that one shot -- not a
CRB or model problem) was enough to dominate the whole-sample RMSE by
itself. Excluding it (a MAD-based robust-outlier filter, $>20\sigma_{\rm
robust}$ in any component; 1/5200 points excluded) brings $\sigma_{vy}$'s
empirical RMSE to $6.77\times10^{-7}$.

\begin{table}[h]
\centering
\begin{tabular}{lccc}
\toprule
Component & CRB (avg $z0$,$z100$) & Empirical RMSE (outlier excl.) & Ratio \\
\midrule
$\mu_{x0}$      & $9.655\times10^{-6}$ & $9.816\times10^{-6}$ & $1.017$ \\
$\mu_{y0}$      & $9.655\times10^{-6}$ & $9.531\times10^{-6}$ & $0.987$ \\
$\mu_{vx0}$     & $2.471\times10^{-6}$ & $2.511\times10^{-6}$ & $1.016$ \\
$\mu_{vy0}$     & $2.473\times10^{-6}$ & $2.437\times10^{-6}$ & $0.985$ \\
$\sigma_x$      & $9.500\times10^{-6}$ & $9.912\times10^{-6}$ & $1.043$ \\
$\sigma_y$      & $9.461\times10^{-6}$ & $9.892\times10^{-6}$ & $1.046$ \\
$\sigma_{vx}$   & $6.767\times10^{-7}$ & $6.861\times10^{-7}$ & $1.014$ \\
$\sigma_{vy}$   & $6.745\times10^{-7}$ & $6.771\times10^{-7}$ & $1.004$ \\
\bottomrule
\end{tabular}
\caption{Full-covariance theta CRB vs.\ empirical RMSE (outlier-excluded),
all 8 kinematic components, confocal\_random\_zernike dataset. All values in
metres (or the matching velocity unit). Every component agrees to within
$\pm5\%$.}
\end{table}

\textbf{Result: the full-covariance CRB matches the empirical scatter to
within a few percent, across all 8 components, once one known bad fit is
excluded.} This directly refutes the presumption (reasonable, given the
prior 66--220$\times$ failure) that Fisher/CRB simply cannot predict
$\theta$'s recoverability -- it was never actually tried with the full
covariance before now.

%%============================================================
\section{Extension: the unaberrated (Gaussian-beam) baselines}
\label{sec:baseline}
%%============================================================

The above was scoped to the aberrated dataset only; this section repeats
\S\ref{sec:comparison}'s CRB-vs-empirical comparison (Step 2+3 only, not the
ridge check) for the two already-fit, unaberrated (\code{CONFOCAL\_FINE})
baseline datasets already saved in \code{non\_phase\_shear/results/}
(\code{kinematic\_estimates\_\{1e6,1e8\}\_N10\_shots50.json}) -- no new data
generation, reusing \code{python-scripts/theta\_crb\_gaussian\_baseline.py}.

\textbf{Data-availability caveat for \code{1e6}.} Its source dataset is no
longer on disk (the only remaining 1e6-atom, 50-shots/run directory has a
different cloud-width prior than what the saved \code{true\_theta} values
actually show, so it isn't the same source). The saved JSON has
\code{true\_theta} and photon-count totals ($n_{\rm tot}$) -- both of which
\code{build\_joint\_fisher} needs -- but not the true $\phi$. Since $\phi$
is drawn independently of $\theta$ (\code{phi0\_mode=random}), this uses a
random $\phi\sim U(0,2\pi)$ substitute per shot instead. \code{1e8} has its
exact source data (\code{data/R40\_N50\_A100000000\_..\_f0.3000}, confirmed
by matching \code{true\_theta}/photon-count against its own file attrs), so
its numbers use the real per-shot $\phi$, no substitution.

\begin{table}[h]
\centering
\begin{tabular}{lcccc}
\toprule
Component & CRB (1e6) & Empirical (1e6) & Ratio & \\
\midrule
$\mu_{x0}$    & $9.644\times10^{-6}$ & $1.017\times10^{-5}$ & $1.054$ & \\
$\mu_{y0}$    & $9.644\times10^{-6}$ & $9.676\times10^{-6}$ & $1.003$ & \\
$\mu_{vx0}$   & $2.500\times10^{-6}$ & $2.638\times10^{-6}$ & $1.055$ & \\
$\mu_{vy0}$   & $2.501\times10^{-6}$ & $2.499\times10^{-6}$ & $0.999$ & \\
$\sigma_x$    & $9.389\times10^{-6}$ & $1.019\times10^{-5}$ & $1.086$ & \\
$\sigma_y$    & $9.367\times10^{-6}$ & $9.990\times10^{-6}$ & $1.067$ & \\
$\sigma_{vx}$ & $6.676\times10^{-7}$ & $7.053\times10^{-7}$ & $1.056$ & \\
$\sigma_{vy}$ & $6.699\times10^{-7}$ & $6.909\times10^{-7}$ & $1.031$ & \\
\bottomrule
\end{tabular}
\caption{\code{1e6} baseline (unaberrated, $\phi$ substituted -- see caveat
above): matches to within $\pm9\%$, similar quality to the aberrated
result.}
\end{table}

\begin{table}[h]
\centering
\begin{tabular}{lcccc}
\toprule
Component & CRB (1e8) & Empirical (1e8) & Ratio & \\
\midrule
$\mu_{x0}$    & $7.616\times10^{-6}$ & $9.279\times10^{-6}$ & $1.218$ & \\
$\mu_{y0}$    & $7.656\times10^{-6}$ & $9.295\times10^{-6}$ & $1.214$ & \\
$\mu_{vx0}$   & $1.972\times10^{-6}$ & $2.402\times10^{-6}$ & $1.218$ & \\
$\mu_{vy0}$   & $1.982\times10^{-6}$ & $2.408\times10^{-6}$ & $1.215$ & \\
$\sigma_x$    & $6.262\times10^{-6}$ & $9.877\times10^{-6}$ & $\mathbf{1.577}$ & \\
$\sigma_y$    & $6.304\times10^{-6}$ & $9.822\times10^{-6}$ & $\mathbf{1.558}$ & \\
$\sigma_{vx}$ & $4.297\times10^{-7}$ & $6.703\times10^{-7}$ & $\mathbf{1.560}$ & \\
$\sigma_{vy}$ & $4.205\times10^{-7}$ & $6.731\times10^{-7}$ & $\mathbf{1.601}$ & \\
\bottomrule
\end{tabular}
\caption{\code{1e8} baseline (unaberrated, real $\phi$, no substitution): a
real, non-trivial gap -- means off by $\sim$21\%, widths off by
$\sim$56--60\%. Confirmed this is not a numerical artifact: the CRB is
stable under $h_\theta$ from $10^{-6}$ to $10^{-8}$ (only degrading at
$10^{-9}$, floating-point roundoff), and the empirical residual distribution
is broadly heavy-tailed (max $\approx4\times$ median, consistent across all
8 components) rather than dominated by one outlier shot as in
\S\ref{sec:comparison}.}
\end{table}

\textbf{Interpretation.} \code{1e6} (both aberrated and unaberrated) and the
aberrated case both match the CRB to a few percent. \code{1e8} does not --
empirical scatter is systematically $1.2$--$1.6\times$ the CRB, worst for
the width parameters. This is not swept under the rug, and a plausible
explanation already exists in prior work rather than needing to be invented
here: \code{notes/joint\_profile\_inference.tex} (\S2026-07-20, a different
optimizer but the same underlying pixel likelihood) documented that
$\theta$'s optimisation \emph{"genuinely struggles at $10^8$"} -- even
warm-started at the exact true $\theta$, a conditional solve "wandered/
plateaued... never approaching tolerance" -- attributed to the likelihood
surface becoming legitimately much sharper at high photon count
($|\nabla_\theta{\rm NLL}|$ exceeding the NLL itself). If
\code{fit\_theta\_best}'s L-BFGS-B is similarly not fully converging at
$10^8$ (not checked here -- see \S\ref{sec:open}), that would inflate the
\emph{empirical} RMSE above the true CRB floor without the CRB itself being
wrong -- consistent with the direction of the observed gap (empirical $>$
CRB, not the reverse) and with $1e6$ (where the same prior work reported
convergence was fine) matching cleanly.

%%============================================================
\section{Framing: this is a single-shot comparison, and that's intentional}
\label{sec:framing}
%%============================================================

$\theta$ is a fresh, independent nuisance parameter drawn from the prior
\emph{every shot} -- it is not a shared, estimated-once parameter like
$\beta=(A_s,A_c)$. Consequently, pooling more shots does \emph{not} reduce
the true per-shot $\theta$ RMSE (there is nothing to pool: shot $i$'s
$\theta_i$ tells you nothing about shot $j$'s $\theta_j$). What more shots
\emph{do} improve is the \emph{measurement precision of the empirical RMSE
estimate itself} -- with 5200 points the empirical RMSE's own relative
uncertainty is $\sim1/\sqrt{2\times5200}\approx1\%$, tight enough to
distinguish the observed few-percent CRB agreement from noise. This is why
this comparison uses many shots (for a trustworthy \emph{measurement}) while
still being a fundamentally single-shot statement (the CRB itself is a
per-shot quantity, not one that shrinks with $N_{\rm shots}$) -- unlike
$\beta$'s CRB, which genuinely does improve with shot count
(\code{crb\_signal.py}'s \code{crb\_from\_j}/\code{simple\_scaling\_crb}).

%%============================================================
\section{From $\theta$-CRB to $\beta$-CRB}
\label{sec:beta}
%%============================================================

\subsection{Method: propagating theta uncertainty into a phase-information number}
\label{sec:beta-method}

\code{crb\_signal.py}'s \code{per\_shot\_information} is exactly the
theta-uncertainty-propagation step: it takes the same joint $17\times17$
Fisher matrix $H$ used in \S\ref{sec:comparison} (over
$\phi,\theta_{z0},\theta_{z100}$) and Schur-complements \emph{only} the 16
theta directions out (not $\phi$), leaving a scalar effective per-shot phase
information $j_{\rm eff}$ -- literally "how much information about the laser
phase $\phi_i$ survives, after accounting for the fact that $\theta$ is
unknown and must be estimated from the same image." Since
$\delta\phi_i(\beta)=A_s\sin(2\pi f i)+A_c\cos(2\pi f i)$ enters only through
$\phi_i$'s argument at $z100$, the $N$-shot Fisher matrix for
$\beta=(A_s,A_c)$ is exactly $F_\beta=\sum_i j_i\,x_ix_i^\top$,
$x_i=[\sin(2\pi fi),\cos(2\pi fi)]$ (\code{crb\_from\_j}), and
${\rm CRB}(\beta)=F_\beta^{-1}$.

\subsection{Converting to rad/$\sqrt{\rm shot}$ and rad/$\sqrt{\rm Hz}$}
\label{sec:units}

Two useful normalisations, both already implemented:

\begin{itemize}
  \item \textbf{rad/$\sqrt{\rm shot}$} (a shot-count-independent number).
    When shots densely sample many signal periods ($\langle x_ix_i^\top
    \rangle\to\tfrac12 I_2$, true here: $f=0.3$ cyc/shot over 50-200 shots is
    15-60 periods), $F_\beta\approx(N j_{\rm bar}/2)I_2$, so
    $\sigma_{A_s}=\sigma_{A_c}\approx\sqrt{2/(N j_{\rm bar})} =
    \sigma_{\rm 1shot}/\sqrt N$ (\code{simple\_scaling\_crb}). $\sigma_{\rm
    1shot}=\sqrt{2/j_{\rm bar}}$ is what "rad/$\sqrt{\rm shot}$" means: it is
    the number such that the CRB at any shot count $N$ is
    $\sigma_{\rm 1shot}/\sqrt N$ -- report this number once, and it answers
    "what precision at $N$ shots" for any $N$ without recomputing anything.
  \item \textbf{rad/$\sqrt{\rm Hz}$} (an amplitude spectral density, the
    standard sensitivity unit in precision metrology). This requires one
    more piece of information this codebase does not define anywhere: the
    real shot \emph{repetition rate} $f_{\rm rep}$ (shots/second) -- every
    frequency in this codebase (\code{--signal\_freq}) is in
    cycles/\emph{shot}, not Hz, so there is no absolute time axis without
    it. Given $f_{\rm rep}$, standard white-noise-ASD conversion applies:
    \[
      {\rm ASD} \;[{\rm rad}/\sqrt{\rm Hz}] \;=\; \sigma_{\rm 1shot}\;[{\rm
      rad}] \times \sqrt{T_{\rm rep}} \;=\; \sigma_{\rm 1shot} /
      \sqrt{f_{\rm rep}},\qquad T_{\rm rep}=1/f_{\rm rep}.
    \]
    (Same logic as converting a per-sample noise std to a PSD/ASD for any
    white, sample-limited process: divide by $\sqrt{\text{sample rate}}$.
    Valid here because each shot's $\theta_i,\phi_i$ are drawn independently
    -- no shot-to-shot correlation to worry about.) $\sigma_{\rm 1shot}$ is
    reported below; multiply by your own $f_{\rm rep}$ whenever you have
    one.
\end{itemize}

Implementation: \code{python-scripts/beta\_crb\_vs\_empirical.py}. Uses a
representative sample of shots (not all $N$) to estimate $j_{\rm bar}$
(checked: relative shot-to-shot spread of $j$ is $\sim$9\% for \code{1e6},
tight enough that treating $j_i\approx j_{\rm bar}$ for all $N$ shots when
building $F_\beta$ via the real $x_i$ pattern is a good approximation, while
avoiding an $O(N)$-shot Fisher-matrix cost).

\subsection{Result: \code{1e6} (Gaussian beam)}
\label{sec:beta-1e6}

\begin{table}[h]
\centering
\begin{tabular}{lccc}
\toprule
& $\sigma_{\rm 1shot}$ [rad] & CRB($N{=}50$) & Empirical ($N{=}50$, \code{best}) \\
\midrule
$A_s$ & \multirow{2}{*}{$1.14\times10^{-3}$} & $1.62\times10^{-4}$ & $4.95\times10^{-4}$ \\
$A_c$ & & $1.62\times10^{-4}$ & $3.21\times10^{-4}$ \\
\bottomrule
\end{tabular}
\caption{$j_{\rm bar}=1.53\times10^6$ (rel.\ shot-to-shot spread $9\%$), 30
sample shots. Ratio (empirical/CRB): $3.05\times$ ($A_s$), $1.99\times$
($A_c$).}
\end{table}

\textbf{Unlike $\theta$'s own CRB (which matched to $\pm5\%$), $\beta$'s CRB
under-predicts the empirical scatter by a real, substantial factor.} Checked
and ruled out one concrete hypothesis before accepting this: rebuilt the CRB
using \emph{exactly} \code{beta\_fits\_from\_kinematics.py}'s own evaluator
geometry (\code{bins=16}, native half-range, vs.\ the tight/32-bin geometry
used elsewhere in this note) -- $j_{\rm bar}$ changed by $<0.3\%$
($1.523\times10^6$ vs.\ $1.526\times10^6$), so evaluator geometry is
\emph{not} the explanation.

\textbf{The actual explanation is structural.} \code{beta\_fits\_from\_kinematics.py}'s
\code{best} method is a \emph{two-stage plug-in} estimator, not the
idealised joint estimator this CRB models:
\begin{enumerate}
  \item Fit $\hat\theta_{\rm best}$ first (a point estimate -- whose own
    scatter matches its CRB, \S\ref{sec:comparison}).
  \item Convert to \code{crb\_signal.py}'s $\eta$ representation via
    \code{theta\_to\_eta}, which \textbf{hardcodes two cross-terms to
    zero}: \verb|[mu_x0, mu_vx0, mu_y0, mu_vy0, sx**2, |\textbf{\texttt{0.0}}\verb|, svx**2, sy**2, |\textbf{\texttt{0.0}}\verb|, svy**2]|
    -- exactly the position-velocity covariance terms.
  \item Fit $\beta$ from that fixed, cross-term-stripped $\eta$ -- $\theta$'s
    uncertainty has already been collapsed to a point estimate and its
    correlation structure discarded, not properly marginalised.
\end{enumerate}
This CRB assumes an efficient estimator that uses the full joint
$(\phi,\theta,\beta)$ information optimally (Fisher-matrix elimination is
the textbook-correct way to get the marginal CRB for \emph{an efficient}
estimator). The real pipeline throws away exactly the correlation structure
$\theta$ was shown to actually have (\S\ref{sec:ridge}'s ridge) partway
through. A two-stage plug-in estimator falling short of the joint bound is a
standard, expected effect, not a bug in either number -- the CRB is
correctly a lower bound; this particular real pipeline doesn't saturate it.

\code{1e8} and the aberrated dataset are pending (see \S\ref{sec:open}) --
this section will be extended once the phase-shear-readout comparison
(also requested, separate note or section) is done.

%%============================================================
\section{Open questions / future work}
\label{sec:open}
%%============================================================

\begin{itemize}
  \item \textbf{Reconcile \S\ref{sec:ridge} and \S\ref{sec:comparison}.}
    Check what the worst-eigenvalue direction $v$ actually projects onto in
    $\theta$-component space, for a few shots -- confirm (or refute) the
    "mixed direction, doesn't dominate any single marginal" hypothesis
    directly rather than leaving it as speculation.
  \item \textbf{Diagnose the $10^8$ gap} (\S\ref{sec:baseline}): check
    whether \code{fit\_theta\_best}'s L-BFGS-B is actually converging
    (\code{res.nit}/\code{res.success}, and/or a tighter \code{maxiter} /
    warm-start-from-truth spot check) at $10^8$ atoms specifically. If
    convergence is the cause, this is an estimator-quality issue, not a
    CRB-validity issue -- worth confirming directly rather than resting on
    the plausible-but-unverified analogy to the joint-solve's documented
    $10^8$ struggles.
  \item \textbf{Exact (not $\phi$-substituted) \code{1e6} baseline} and a
    dedicated apples-to-apples aberrated-vs-Gaussian comparison at matched
    $N$: the \code{1e6}/\code{1e8} baselines used here are convenient
    (already fit, no new compute) but neither is parameter-matched to
    \code{confocal\_random\_zernike} as closely as
    \code{data/R80\_N200\_A1000000\_..\_f0.3000} would be (exact $n_{\rm
    atoms}$/cloud-prior/shot-count match, already on disk, just never run
    through \code{generate\_kinematic\_estimates.py}) -- worth doing if a
    tight, controlled aberration-only comparison is needed later.
  \item \textbf{Connect to the $\beta$ (signal) CRB.} Propagate whatever
    aberrated-vs-Gaussian theta-CRB difference is established through
    \code{crb\_signal.py}'s existing $j_{\rm eff}\to$ beta-CRB machinery to
    get an actual quantitative bound on how much wavefront
    aberration/curvature costs in signal-extraction precision -- the
    original motivating question for this whole line of work.
  \item The current check used only \code{run\_000}'s first 10 shots for the
    CRB side (cheap: $\sim1$ minute for all of \S\ref{sec:ridge}--\ref{sec:comparison}
    combined). Widening this to more shots/runs would tighten the CRB-side
    estimate too (currently it isn't the limiting source of noise in the
    comparison, but worth checking as the empirical run finishes all 20
    runs).
\end{itemize}

\end{document}
